% Retrospective section. No preamble; requires amsmath, amssymb and url/hyperref.
% Mathematical records below concern original W, not its ordinary order extension.
\section{Cardinal order, exponentiation, and choice in the original model}
\label{sec:complete-cardinals}

This section collects the results established for the \emph{original}
symmetric model $W$, and, where explicitly indicated, its separately audited
parameter variants. It does not transfer these conclusions to the later
ordinary forcing extension. The foundational construction, its Partition
Principle and support lemmas are due to OpenAI~\cite{OpenAI}; the additional
calculations below are the research deductions reported here. They have
received separate mathematical checks by AI agents, not external human peer
review or a complete Lean formalization. Historical qualifications are given
at the end of the section.

\subsection{The generator and the exact cardinal records}

Write $X\leq Y$ for an injection and $X\cong Y$ for a bijection. In the
source's general ground $V\models\mathsf{ZFC}$, put
\[
 \kappa=\omega_1^V,\qquad \mu=((2^{\aleph_0})^+)^V.
\]
The set $A\subseteq\mathcal P(\kappa)$ consists of the distinct Cohen rows.
The source proves $\mathsf{SVC}(A)$, choice for ordinal-indexed families
$\mathsf{AC}_{\mathrm{WO}}$, and the classification
\[
 \text{every non-well-orderable image of $A$ is bijective with $A$}. \tag{Q}
\]
Together these give $\mathsf{PP}$ and $\neg\mathsf{AC}$ in $W$.
The same source locality implies that a function on an ordinal with values
in a ground set belongs to the Cohen intermediate extension. Thus ground
cardinals at least $\mu$ remain cardinals in $W$.

\paragraph{The action and absorption.}
The generic diagrams define an internal free right action of the ground
monoid $T$ on $A$. For each $a$, the map $t\mapsto a\cdot t$ is injective;
its image $aT$ is a \emph{chart}, of cardinal $\mu$. Every ordinal-indexed
subset of $A$ lies in one chart. The latter assertion follows by copying a
putative outside value freshly while fixing the supporting diagram: two
compatible conditions would assign different labels to the same ordinal
index. Hence the largest well-orderable cardinal injecting into $A$ is
\[
                         w(A)=\mu,\qquad h(A)=\mu^+.
\]
Here $h$ is the Hartogs number. No global inverse-label map $A\to I$ is
asserted to belong to $W$.

Each map $a\mapsto a\cdot t$ is onto. Define a predecessor diagram by
$q(tx)=p(x)$ and put its remaining coordinates freshly; exact amalgamation
makes such predecessors dense. For $t\ne1$, every fiber over an anchor
$a_0$ is non-well-orderable. To see the key step, normalize an alleged
ordinal enumeration to a condition $(d,p)$ containing its support and the
Cohen support, with $a_0=A_{p(ce_0)}$, where $e_0\ne1$ is the distinguished
monoid element. Take $q_{e_0}=p$. The source probe clause at
$(e_0c,t,\{1\})$ supplies $n$ with
\[
                   nt=e_0ce_0,\qquad e_0\not\sqsubseteq n,
\]
where $u\sqsubseteq v$ means $uz=v$ for some $z$.
Thus $A_{q(n)}$ is in the fiber but outside the root diagram. Decide its
enumeration index and copy both the stronger diagram and its Cohen
condition freshly over the root. Exact diagram overlap and disjoint
Cohen supports off the root give the contradiction.

Fibers for distinct $t$ are disjoint by freeness in the monoid coordinate.
Each nonempty fiber is an image of $A$ by a retraction, so (Q) identifies
it with $A$. Choosing these bijections over the well-ordered set
$T\setminus\{1\}$ gives $A\times\mu\leq A$. Consequently
\[
                  A\times\mu\cong A,\qquad A^2\cong A.       \tag{1}
\]
For the second equality, first use the raw normal form below for $A^2$.
Its well-orderable subsets have well-orderable projections of size at most
$\mu$, so both ordinal parameters are at most $\mu$ and are absorbed by
$A\times\mu$. This avoids assuming square absorption in advance.

\paragraph{Normal form and order.}
An $\mathsf{SVC}(A)$ presentation $A\times\eta\twoheadrightarrow X$ splits
$X$ by its least presenting ordinal coordinate. By (Q), each resulting
non-well-orderable piece is a copy of $A$; the other pieces are
well-orderable. Ordinal-indexed choice assembles the bijections, yielding
$X\cong(A\times\alpha)\sqcup\beta$.
Using (1), the exact records are
\[
 \begin{array}{ll}
 (0,\lambda):&\text{a well-orderable cardinal }\lambda,\\
 (a,e),\ \mu\leq a\leq e:
   & A\times a\quad\text{if }a=e,\\[-2pt]
   &(A\times a)\sqcup e\quad\text{if }a<e.
 \end{array}                                                    \tag{2}
\]
All nonzero coordinates here are ordinal cardinals. For the second kind,
$w(X)=e$, $h(X)=e^+$, and
\[
             A\leq X,\qquad e\leq X\leq A\times e.
\]
Under $\mathsf{AC}_{\mathrm{WO}}$, the Lindenbaum number agrees with the
Hartogs number. The least non-well-orderable cardinal is therefore $|A|$.

The records determine the \emph{entire injection order}:
\[
            (a,e)\leq(c,f)\quad\Longleftrightarrow\quad
                         a\leq c\ \hbox{and}\ e\leq f,        \tag{3}
\]
including the well-orderable records with first coordinate zero.
For necessity in the first coordinate, an injection between non-WO
normal forms cannot send an entire source $A$-slice into the target's
ordinal tail. Ordinal-indexed choice picks one point from each source
slice whose image lies in the target's $A$-part. These are distinct,
so the source coefficient is at most the largest WO cardinal of that
part, namely the target coefficient. Necessity for $w$ is immediate.
Conversely embed the $A$-part and the ordinal part using the two coordinate
inequalities, then combine the embeddings using doubling of the target.
The WO cases reduce to the definition of $w$.

\subsection{Lattice structure, finite bases, and finite powers}

Equation (3) gives a precise cardinal lattice, with necessary size
qualifications: \emph{every set-sized family} of cardinalities has a
supremum, and every \emph{nonempty set-sized family} has an infimum.
They are obtained by coordinatewise suprema and minima in (2).
The resulting pairs are again realizable records. Binary distributive
laws follow from those for minimum and maximum in two ordinal chains.
There is no largest cardinal and no empty infimum; this is not an
unqualified claim of a complete lattice or of complete distributivity.
Scott cardinals can be used throughout to avoid choosing representatives.

The finite-basis theorem proved elsewhere in this report says that every
nonempty family has a finite coinitial family under injections. It yields
$\mathsf{NDS}$ and the absence of infinite antichains of cardinalities.
It does \emph{not} bound all antichains by one fixed natural number. For
any positive $n$, the $n$ records
\[
 \bigl(\mu^{+i},\mu^{+(2n-i)}\bigr),\qquad i<n,
\]
form an antichain: the first coordinates increase and the second decrease.
In bare ZF the exact order-theoretic equivalence is
$\mathsf{FB}\leftrightarrow\mathsf{WF}_{\mathrm{card}}+\mathsf{FAC}$,
where $\mathsf{WF}_{\mathrm{card}}$ is minimal-element well-foundedness
on sets of Scott cardinals and $\mathsf{FAC}$ means every such antichain
is finite. With DC it becomes
$\mathsf{FB}\leftrightarrow\mathsf{NDS}+\mathsf{FAC}$.
The representative and sequence distinctions should not be suppressed
when DC is absent. Moreover, the already proved PP and NDS give all
24 cardinal well-foundedness variants in Blass--Kulshreshtha's
Definition~4.1, using their equivalence under PP
\cite[Proposition 9.8]{BK}. This equivalence does not prove NDS from
PP alone.

For a fixed ordinal cardinal $\lambda$, the assertion
$\mathsf W_\lambda$ that every cardinal is comparable with $\lambda$
has the exact threshold
\[
                W\models\mathsf W_\lambda
                       \quad\Longleftrightarrow\quad\lambda\leq\mu.
                                                                  \tag{4}
\]
Every non-WO set contains $A$, proving the positive direction; $A$ is
incomparable with every ordinal cardinal above $\mu$, proving the negative.
The usual surjective version agrees under PP. In particular full
comparability still fails, as it must when AC fails.

Every infinite set absorbs a countable multiplier and doubles. For
non-WO $X$ with record $(a,e)$, direct multiplication gives
\[
                 X^n\cong A\times e\quad(2\leq n<\omega).     \tag{5}
\]
Indeed $A\times e\leq X^2\leq A^2\times e^2$, and multiplying the
normal form by $A\times e$ leaves the same cardinal. Thus
$X^3\cong X^2$ holds for every infinite $X$, but $X^2\cong X$ does not:
$X=A\sqcup\mu^+$ has record $(\mu,\mu^+)$ and its square has record
$(\mu^+,\mu^+)$. The question about universal square/cube stability is
historical; Rubin--Rubin's 1985 introduction already discusses it. The
bounded search does not establish its intervening publication status.

There is also a finite bound, depending on $X$, on iterated power-set
roots: some $N<\omega$ satisfies
\[
                    \mathcal P^n(Y)\cong X\quad\Longrightarrow\quad n\leq N.
                                                                  \tag{6}
\]
Represent every possible root by a subset of $X$, using iterated singleton
maps. The resulting set of roots has a finite coinitial basis
$Y_1,\ldots,Y_k$, with unique depths $n_1,\ldots,n_k$; uniqueness follows
from Cantor's theorem. If a root $Y$ had depth $m>\max n_i$, some
$Y_i\leq Y$ would give
$X\cong\mathcal P^{n_i}(Y_i)\leq\mathcal P^{n_i}(Y)
 <\mathcal P^m(Y)\cong X$, impossible. In particular every cardinal is
a finite power-set iterate of one which is not a power-set cardinal.
There is no bound uniform in $X$.

\subsection{Supported names and exact exponentiation}

For ground ordinals $b,\lambda$, write
\[
             E_V(b,\lambda)=|\{f\in V:f:\lambda\to b\}|^V.
\]
This is a ground cardinal; unmarked exponentials are computed in $W$.
The decisive local calculation, for $b\geq2$ and ordinal cardinals
$\lambda\geq\mu$, is
\[
 b^A\cong A\times E_V(b,\mu),\qquad
 b^\lambda\cong A\times E_V(b,\lambda).                         \tag{7}
\]
The coefficients are at least $(2^\mu)^V$ and are preserved cardinals.

Here are the ingredients ensuring that (7) counts \emph{all internal}
well-orderable families. Fix a supporting diagram $p$. Normalize names
to globally total $A$-to-$b$ functions, using the zero function whenever
the original value has the wrong type. The canonical graph name uses
only label/ordinal pair names and the ordinary forcing relation. It is
equivariant, keeps the support, and has a fixed rank bound. Thus these
names form a ground set, rather than a class of unbounded-rank names.
For each $t\in T$ choose a ground diagram $q^t$ with $(q^t)_t=p$.
At the record $(t,s,d)$, with $d$ a Cohen condition on $D_{q^t}$, record
$\xi<b$ exactly when
\[
                 (d,q^t)\Vdash\tau(A_{q^t(s)})=\xi.
\]
There are $\mu$ positions, and at most one value per position. The code
is a \emph{partial value function}, so there are at most
$(|b|^V+1)^\mu=E_V(b,\mu)$ codes in the ground.

Equal codes force equality below $(1,p)$. Otherwise put an input where
the values differ into a larger diagram $r$, write $p=r_t$, restrict
the Cohen witness to $D_r$ by source locality, and transport it to $q^t$
by a permutation fixing $D_p$. It produces different values at the same
record, a contradiction. Ground choice now selects one representative
per code. Their ordinal enumeration is HS with support $D_p$.
Any internal WO family of functions has one supported enumeration;
normalizing its coordinates with the same support embeds the whole family
into that representative list by least indices. Ground functions on
one chart give the matching lower bound. For ordinal domain $\lambda$
the records are $(d,\xi)$, $\xi<\lambda$, and both local-decision clauses
give the analogous count. There is no maximum-principle assumption,
choice indexed by $A$, or automorphism of the actual pure model in this
argument. Finally both function sets are non-WO and square-idempotent;
(2)--(5) turn their exact WO bounds into (7).

\paragraph{Universal exponent collapse.}
For every non-WO $X$, set $e=w(X)$. Then for \emph{every set} $Y$,
\[
                         Y^X\cong Y^e.                         \tag{8}
\]
More explicitly,
\[
 \begin{aligned}
 \mathcal P(X)&\cong A\times E_V(2,e)\cong\mathcal P(e),\\
 b^X&\cong A\times E_V(b,e) &&(b\geq2\text{ ordinal}),\\
 Y^X&\cong A\times E_V(f,e) &&(Y\text{ non-WO},\ f=w(Y)).
 \end{aligned}                                                   \tag{9}
\]
The same formulas hold with domain $e$.

For the proof, two-point coding and $A\leq\mathcal P(\kappa)$ give
$A^e\cong\mathcal P(e)$. The normal form also gives
$\kappa\times X\cong X$, hence $A^X\cong\mathcal P(X)$.
Sandwich an ordinal-base power between $b^e$ and $b^{A\times e}$.
With $\delta=E_V(b,\mu)$, (7) calculates the upper endpoint as
\[
 (b^A)^e\cong(A\times\delta)^e
 \cong A^e\times\delta^e\cong A\times E_V(b,e),
\]
since ground arithmetic has $E_V(\delta,e)=E_V(b,e)$.
The lower endpoint has the same cardinal. Its case $b=2$ proves the
power-set formula before that formula is used for non-WO bases.
For non-WO $Y$, use $A\leq Y$, $f\leq Y\leq A\times f$ for the upper
bound. To inject $A\times(f^e)^V$ into $Y^e$, reserve one coordinate for
the $A$ point and $e$ further coordinates for the ground function,
using $e+1\cong e$. This provides the matching lower bound in both
$Y^e$ and $Y^X$. Empty and singleton bases give the trivial equalities
because $X$ and $e$ are nonempty. This exhausts all internal bases.

Consequently
\[
                  \mu<A,\qquad \forall Y\ (Y^\mu\cong Y^A).    \tag{10}
\]
The statement asserts pointwise existence of bijections. It asserts
neither a natural transformation between function-space functors nor
one class-sized choice of all bijections.

\subsection{Complete arithmetic table over the constructible ground}

For the separately audited $L$-ground variants $M_{m,n}$, put
$\kappa=\aleph_m$, $\mu=\aleph_n$, where $1\leq m<n<\omega$.
All alephs are preserved. The original construction over $L$ is $M_{1,2}$.
The following tables concern these original symmetric models.
The records of power sets are
\[
\begin{array}{c|c}
 |X|&|\mathcal P(X)|\text{ as a record}\\ \hline
 \text{finite }r&(0,2^r)\\
 \text{infinite WO }\lambda<\kappa&(0,\lambda^+)\\
 \text{WO }\kappa\leq\lambda<\mu&(\mu,\mu)\\
 \text{WO }\lambda\geq\mu&(\lambda^+,\lambda^+)\\
 \text{non-WO }(a,e)&(e^+,e^+)
\end{array}                                                       \tag{11}
\]
For the small ordinal plateau, the local Cohen forcing has size $\mu$
and the $\kappa^+$ chain condition. Nice names for subsets of $\lambda$
number at most $\mu^{\max(\kappa,\lambda)}$, equal under ground GCH to
$\mu$ for $\kappa\leq\lambda<\mu$ and to $\lambda^+$ for
$\lambda\geq\mu$. One supported enumeration bounds every WO family;
a chart and the ground power set supply the lower bounds. Below $\kappa$
there are no new subsets of ordinals. In particular, in $M_{1,2}$,
\[
 \mathbb R\cong\aleph_1,\qquad
 \mathcal P(\aleph_1)\cong A,\qquad
 \mathcal P(A)\cong\mathcal P(\aleph_2)\cong A\times\aleph_3.
\]
This is a ground-GCH calculation, not GCH inside the choiceless model.

For completeness write $E(b,\lambda)=E_L(b,\lambda)$. Apart from empty
and singleton cases, the \emph{entire exponent table} is
\[
\begin{array}{c|c|c}
 |Y|&|X|&|Y^X|\text{ as a record}\\ \hline
 \text{WO }b\geq2&\text{WO }2\leq\lambda<\kappa&(0,E(b,\lambda))\\
 \text{WO }b\geq2&\text{WO }\lambda\geq\kappa&(\theta,\theta)\\
 \text{non-WO }(a,f)&\text{WO }\lambda\geq2&(E(f,\lambda),E(f,\lambda))\\
 \text{WO }b\geq2&\text{non-WO }(c,e)&(E(b,e),E(b,e))\\
 \text{non-WO }(a,f)&\text{non-WO }(c,e)&(E(f,e),E(f,e))
\end{array}                                                       \tag{12}
\]
where $\theta=\max(\mu,E(b,\lambda))$.
For $X=\varnothing$ the result is $1$, including $0^0=1$;
for $X\ne\varnothing$ and $Y=\varnothing$ it is $0$;
a singleton base always gives $1$; a singleton exponent returns $Y$
with its original record. This last exception matters when $a<f$.

To justify the rows not already covered by (9), ordinal-valued functions
of length $\lambda<\kappa$ are ground functions. Also
$A^\lambda\cong A$ for $0<\lambda<\kappa$: a sequence has range in one
chart and its $T$-coordinate sequence is ground, giving
$A\times(T^\lambda)^V\twoheadrightarrow A^\lambda$; now
$\mu^{<\kappa}=\mu$, (1), and PP apply. For $\lambda\geq\kappa$,
nice function names number at most
$(\max(\mu,b)^\lambda)^L=\max(\mu,E(b,\lambda))$.
This gives the WO-base row as above. The non-WO-base upper bound is
$A^\lambda\times f^\lambda$. A matching lower bound uses two coordinates
if $\lambda$ is finite, since $f^\lambda=f$ in the ground, and the
reserved-coordinate coding if $\lambda$ is infinite.

The ground coefficients in (12) are explicit. For positive finite
$\lambda$, $E(b,\lambda)$ is the usual finite power when $b$ is finite
and is $b$ when $b$ is infinite. For infinite $\lambda$ and $b\geq2$,
\[
 E(b,\lambda)=
 \begin{cases}
  \lambda^+,&b\leq\lambda,\\
  b,&\lambda<b\text{ and }\lambda<\operatorname{cf}^L(b),\\
  b^+,&\operatorname{cf}^L(b)\leq\lambda<b.
 \end{cases}                                                     \tag{13}
\]
The final case concerns singular $b$; it follows from K\"onig's theorem
and $b^\lambda\leq2^b=b^+$.

\subsection{Choice consequences and their exact boundaries}

In addition to $\mathsf{PP}$, $\mathsf{AC}_{\mathrm{WO}}$,
$\mathsf{FB}$ and $\mathsf{NDS}$, the original model satisfies ordinary
DC, $\mathsf{KWP}_1$, the Ordering Principle, and choice for arbitrary
families of nonempty finite sets. The classical implication
$\mathsf{AC}_{\mathrm{WO}}\Rightarrow\mathsf{DC}$ is Jensen's theorem
as presented in Jech~\cite[Theorem 8.2]{Jech}; an earlier apparent source
conflict was resolved by reading the actual PDF. For $\mathsf{KWP}_1$,
PP turns a presentation into $X\leq A\times\eta$, which embeds into a
power set of an ordinal by tagged coding. Lexicographic order on that
power set linearly orders $X$. Least elements of finite subsets then
give finite-member choice. This does not assert that every prescribed
partial order has a linear extension: the separate obstruction in this
report proves failure of that stronger principle and of BPI.

The source itself explicitly states the dual Cantor--Bernstein principle
$\mathsf{CSB}^*$ and Weak Partition Principle as consequences of PP.
With $X\leq^*Y$ meaning $X=\varnothing$ or $Y\twoheadrightarrow X$,
these say respectively
\[
 X\leq^*Y\ \&\ Y\leq^*X\Rightarrow X\cong Y,
 \qquad X\leq^*Y\Rightarrow\neg(Y<X).
\]
They are not newly extracted results of this report.

\paragraph{Choice for countable members.}
A further saved proof, now independently audited twice, gives
\[
 \mathsf C(\infty,\leq\aleph_0):\quad
 \text{every family of nonempty countable sets has a choice function}.
                                                                  \tag{14}
\]
This is Howard--Rubin Form 85A. The family of indices is arbitrary;
its \emph{members} are countable, so (14) differs from countable choice.
For one nonempty countable $S\subseteq A\subseteq\mathcal P(\omega_1)$,
its pairwise first-difference ordinals are countably many and therefore
bounded. Let $\alpha(S)<\omega_1$ be the least ordinal whose restriction
map separates $S$. Source preservation gives
$\mathcal P(\alpha)^W=\mathcal P(\alpha)^V$ for every $\alpha<\omega_1$.
One fixed ground family of well-orders of these power sets selects the
least restriction and hence a unique member of $S$. No simultaneous
choice of enumerations of $S$ was used. Given arbitrary $X$, choose one
embedding $X\leq A\times\eta$ and first take the least ordinal slice of
a countable subset, then this selector on $A$. For an arbitrary family,
apply the construction just once to its union. The proof is for original
$W$; new reals prevent an automatic transfer to the later extension.

\paragraph{Ordinary versus higher dependent choice.}
Despite DC, the original model and the audited parameter family fail
$\mathsf{DC}_{\omega_1}$. Order $A$ natively by
$a\leq_Tb$ iff $a\cdot t=b$ for some $t\in T$. Every ordinal-indexed
subset has a common lower bound by chart locality, and every point has
a strict predecessor by surjectivity of the $e_0$-translation.
There is no strictly decreasing $\omega_1$-sequence: localize its whole
range to one chart, lift to $T$, and apply the source's strictly
increasing real-valued height. This would give a decreasing
$\omega_1$-sequence of reals, impossible because its disjoint successive
intervals inject into the rationals. Thus the tree of descending
sequences of countable ordinal length is countably closed and has no
terminal node but has no $\omega_1$-branch. Equivalently, the relation
``strictly below every member of the preceding countable set'' directly
violates the conclusion of Howard--Rubin's Form 44.
Increasing only the Cohen parameter does not remove this height obstruction.

\subsection{Ramsey principles: ordinary truth and stronger separations}

Here RPP means the \emph{Ramseyan Partition Principle}: every two-coloring
of pairs of an infinite set has an infinite homogeneous subset. This is
not PP and not the group-filter Ramsey condition used in FM analyses of
BPI. Ordinary RPP holds in $W$: DC gives a countably infinite subset,
and the ordinary finite-arity Ramsey theorem on $\omega$ is provable in
ZF. More generally, $\mathsf{NDS}$ implies that every infinite set is
Dedekind-infinite: otherwise finite injective lists from an infinite
Dedekind-finite set produce a strictly descending sequence by excluding
successively shorter lists. Thus both the PP/DC and NDS routes give RPP.
These implications are classical weak-choice context~\cite{Blass1977,BK}.

The stronger cardinal demands differ. Let $\mathrm{NR}_\lambda$ say
that every two-coloring of pairs of every non-WO set has a homogeneous
subset of cardinality $\lambda$. Throughout the audited parameter family,
\[
                 \mathrm{NR}_\lambda
                          \quad\Longleftrightarrow\quad\lambda\leq\kappa
                 \quad(\lambda\text{ infinite WO}).            \tag{15}
\]
For the positive part pull a coloring back to a $\kappa^+$-subset of
$A$. An inner $L[C]$ containing its ordinal code and the required small
sequence enumerations satisfies the appropriate Erd\H{o}s--Rado theorem;
its witness remains in $W$. The first-difference parity coloring on
$A\subseteq\mathcal P(\kappa)$ gives the negative part: fresh-copy
arguments make every homogeneous set WO, and a Cohen-condition
$\Delta$-system argument bounds its cardinal by $\kappa$. Each color
also has a homogeneous set of size $\kappa$. The arbitrary-ground pair
proof uses the separately audited collapse of $2^{<\kappa}$ to $\kappa$
in $W$; the $L$-ground cases need only their unchanged smaller power sets.

The color-count boundary is also exact. For every ordinal cardinal
$1\leq\eta<\kappa$, every $\eta$-coloring of pairs of a non-WO set has a
$\kappa$-sized homogeneous subset. At $\eta=\kappa$ the coloring
$\{a,b\}\mapsto\min(a\mathbin\triangle b)$ has no homogeneous triple:
three distinct binary bits cannot be pairwise different at one coordinate.
Thus $\kappa$ is the least WO number of colors permitting such a failure.

For the $L$-ground family $M_{m,n}$, let
$\mathrm{NR}^{r,\mathrm{fin}}_\lambda$ assert that every coloring of
$[X]^r$ into any positive finite number of colors, for every non-WO $X$,
has a homogeneous set of cardinality $\lambda$. Its exact spectrum is
\[
 M_{m,n}\models\mathrm{NR}^{r,\mathrm{fin}}_\lambda
 \quad\Longleftrightarrow\quad
 \lambda\leq\aleph_{\max(m-r+2,0)},\qquad r\geq2.             \tag{16}
\]
Here $\lambda$ ranges over infinite WO cardinals. Inner-model Erd\H{o}s--Rado
supplies the positive bound; a generalized Sierpi\'nski coloring and
first-difference stepping up give the negative bound~\cite{ErdosRado,BHT}.
The proof checks that homogeneous sets in the stepping-up construction
are well-ordered by lexicographic order or its reverse, rather than
assuming this in advance. The triple obstruction uses \emph{four colors}.
The exact two-color triple question remains unresolved here; the audit
of simple binary coarsenings does not settle it.

For example, $M_{1,3}$ and $M_{2,3}$ have the same
$w(A)=\aleph_3$, $h(A)=\aleph_4$, and the same entire threshold (4).
Their pair homogeneous thresholds are $\aleph_1$ and $\aleph_2$;
their finite-color triple thresholds are $\aleph_0$ and $\aleph_1$.
Both retain PP, FB, NDS, $\mathsf{AC}_{\mathrm{WO}}$, DC and failure of
$\mathsf{DC}_{\omega_1}$ and AC. This separates the stronger Ramsey
statements within that shared package. It is not a claim that the
classical colorings or partition calculus are new.

\subsection{Comparability, cancellation, and historical boundaries}

All power sets in $W$ have comparable cardinalities. An infinite set
doubles, so its power set squares to itself. Every non-WO power set
therefore has record $(\theta,\theta)$, and these form a chain. A WO
power set must have a WO base of cardinal below $\kappa$; otherwise it
contains $A\leq\mathcal P(\kappa)$. Such smaller powers have ground
enumerations below $\mu$, so lie below every non-WO power. Thus
\[
 \mathsf{PC}:\qquad
  \forall X,Y\ [\mathcal P(X)\leq\mathcal P(Y)
                  \ \text{or}\ \mathcal P(Y)\leq\mathcal P(X)]. \tag{17}
\]
Bare $\mathsf{PC}+\neg\mathsf{AC}$ is prior art: Glazer's publicly dated
30 September 2025 answer proves it in Cohen's first model~\cite{Glazer}.
The joint PP/FB/DC package and exact coefficients are the application
under discussion here; no exhaustive novelty guarantee is claimed.
In ZF, PC also implies $\mathsf{KWP}_1$: compare $\mathcal P(X)$
with $\mathcal P(\theta)$ for $\theta=h(\mathcal P(X))$. The latter
cannot inject into the former, so $X\leq\mathcal P(X)\leq\mathcal P(\theta)$.
This observation appears in Wojowu's comment on the cited question.

The Injective Continuum Function assertion
$\mathsf{ICF}:\mathcal P(X)\cong\mathcal P(Y)\Rightarrow X\cong Y$
fails already in ZFC at a permissible continuum plateau, and is not a
consequence of AC. Here it fails more specifically at the strict pair
$\mu<A$, by (10). The order-reflecting variant
$\mathcal P(X)\leq\mathcal P(Y)\Rightarrow X\leq Y$ fails as well.
For finite Kinna--Wagner principles,
\[
          \mathsf{ICF}+\mathsf{KWP}_{n+1}
                           \Rightarrow\mathsf{KWP}_{n}\quad(n<\omega),
                                                                  \tag{18}
\]
so ICF plus any finite $\mathsf{KWP}_n$ implies AC.
Apply $\mathsf{KWP}_{n+1}$ to $\mathcal P(X)$ to get
$\mathcal P(X)\leq\mathcal P(B)$ with $B=\mathcal P^n(\theta)$,
$\theta$ infinite. Since $B$ doubles,
$\mathcal P(X\sqcup B)\cong\mathcal P(B)$; ICF gives $X\leq B$.
Here $\mathsf{KWP}_0=\mathsf{AC}$. No transfinite version or priority
claim is asserted~\cite{KS,ICFquestion}. The corresponding finite
surjective principles $\mathsf{KWP}^{*}_n$ imply $\mathsf{KWP}_{n+1}$
by taking inverse images under the presentation surjection. Thus ICF
combined with any finite $\mathsf{KWP}^{*}_n$ implies AC as well.

For clarity, distinguish two all-base assertions: (EC) if
$Y^X\cong Y^Z$ for every $Y$, then $X\cong Z$; and (SES) if $X<Z$,
some $Y$ satisfies $Y^X<Y^Z$. These are descriptive labels, not asserted
standard names. Both hold in ZFC. For infinite cardinals $\alpha<\beta$,
choose $\nu=\beth_{\alpha^+}$, a strong limit of cofinality $\alpha^+$;
then $\nu^\alpha=\nu<\nu^\beta$ by bounded ranges and K\"onig's theorem.
Finite cases use base $2$. Equation (10) refutes both in $W$.

Bare comparable all-base plateaux have an older mechanism. Mutual
surjections induce injections between function spaces by precomposition
at every base. From an infinite Dedekind-finite set $D$, let $S$ be its
finite injective lists. Deleting the first entry, with the empty list
sent to an extra point, maps $S$ onto $S\sqcup1$; a retraction maps back.
But $S<S\sqcup1$, so their powers agree at every base. Hamkins's 2012
answer records this list construction; the all-base conclusion is the
immediate precomposition extension~\cite{Hamkins}. Thus bare failure of
EC or SES is not a new consistency phenomenon. The distinction in (10)
is its compatibility with PP, DC and dual CSB, and its one well-orderable
exponent; it cannot arise from unequal mutually surjective sets.
For the logical record, ZF proves
$\mathrm{EC}\Rightarrow\mathrm{SES}$: on a comparable nonempty pair,
function extension gives the forward injections for every nonempty base;
if none were strict, EC would give equality of the exponents. An empty
smaller exponent is separated by base $2$. It also proves
$\mathrm{EC}\Rightarrow\mathsf{CSB}^*$, by the precomposition argument,
and either EC or SES implies that every infinite set is
Dedekind-infinite, by the finite-list counterexample just given.
No converse or general equivalence with AC is established here.
Finally strict \emph{reflection},
$p^x<p^y$ and $p\ne0\Rightarrow x<y$, is a different statement, listed
as equivalent to AC by Lindenbaum--Tarski~\cite[No.~82(A5)]{LT}.
Its direction and quantifiers do not identify the strength of EC or SES.

The general implication $\mathsf{PP}\Rightarrow\mathsf{NDS}$ has not been
proved by these model calculations. Nor do they prove PP alone implies
FB, the exact arithmetic, or the stronger Ramsey statements. Earlier
phase notes that left BPI in this concrete model undecided are superseded
by the separate negative result in this report. Positive BPI in a different
PP model, and the global choice classification of the later extension,
remain distinct questions.

\paragraph{Proof-file map and verification scope.}
Detailed records are in \nolinkurl{evidence/original_model/}:
\nolinkurl{T_ACTION_FIBERS_AND_SQUARE_COEFFICIENT_DRAFT.txt};
\nolinkurl{T_ACTION_INDEPENDENT_AUDIT_AND_DC_OMEGA1.txt};
\nolinkurl{EXACT_EXPONENTIALS_AND_UNIVERSAL_EXPONENT_COLLAPSE_AUDIT.txt};
\nolinkurl{COMPLETE_GCH_GROUND_EXPONENT_TABLE.txt};
\nolinkurl{FB_NDS_FAC_EXACT_CHOICE_AUDIT.txt};
\nolinkurl{COUNTABLE_MEMBER_CHOICE_SECOND_INDEPENDENT_AUDIT.txt};
\nolinkurl{GLOBAL_NON_WELL_ORDERABLE_RAMSEY_THRESHOLD.txt};
\nolinkurl{HIGHER_ARITY_ROOT_INDEPENDENT_AUDIT.txt}; and
\nolinkurl{ALL_BASE_EXPONENT_CANCELLATION_LITERATURE_AUDIT.txt}.
The lattice and root-depth proofs are in
\nolinkurl{evidence/finite_basis/PROOFS_cardinal_lattice_and_power_roots.txt}.
These files preserve source dependencies and search limits. The new results
in this section are mathematical deductions, not a certification of
novelty or a claim that all of them have been formally verified.
